% J1.4: tres formas de escribir la misma cadena de operaciones.
\documentclass[tikz,border=4pt]{standalone}
\usepackage{resonancia}
\begin{document}
\begin{tikzpicture}
  \node[etiqueta, anchor=west] at (0,1.9) {\textbf{Anidada}};
  \node[anchor=west, font=\ttfamily\small] at (2.2,1.9) {rms(quitar\_media(x))};
  \node[etiqueta, anchor=west] at (0,1.1) {\textbf{Con tubería}};
  \node[anchor=west, font=\ttfamily\small] at (2.2,1.1) {x |> quitar\_media |> rms};
  \node[etiqueta, anchor=west] at (0,0.3) {\textbf{Con composición}};
  \node[anchor=west, font=\ttfamily\small] at (2.2,0.3) {(rms $\circ$ quitar\_media)(x)};
  % la cadena
  \node[cajaAzul, minimum width=1.2cm] (x) at (1.0,-1.3) {x};
  \node[cajaVerde, minimum width=2.6cm] (q) at (4.0,-1.3) {quitar\_media};
  \node[cajaNaranja, minimum width=1.6cm] (r) at (6.9,-1.3) {rms};
  \node[cajaAzul, minimum width=1.4cm] (y) at (9.0,-1.3) {0,773};
  \draw[flecha] (x) -- (q);
  \draw[flecha] (q) -- (r);
  \draw[flecha] (r) -- (y);
  \node[etiqueta, color=rGris, anchor=north] at (5.0,-2.0) {la misma operación, de izquierda a derecha};
\end{tikzpicture}
\end{document}
